%! Function Model Construction and Application — Unit 1, Lesson 1.14 — Warm-Up
\documentclass[10pt]{article}
\usepackage{precalculus-article}
\usepackage{precalculus-boxes}
\newcommand{\TallMath}[1]{$\displaystyle #1\rule[-1.4em]{0pt}{3.2em}$}

\begin{document}

\pageheader{Unit 1, Lesson 1.14}{Warm-Up}
\namedateperiod

\vspace{0.12in}

\begin{enumerate}[itemsep=18pt]

\item A table of values for $f(x)$ is given below (where $x$ is in seconds and $f(x)$ is in feet).

      \smallskip
      \renewcommand{\arraystretch}{1.3}
      \begin{tabular}{|c|c|c|}
      \hline
      $x$ & 2 & 10 \\
      \hline
      $f(x)$ & 8 & 32 \\
      \hline
      \end{tabular}
      \smallskip

      Compute the average rate of change of $f$ on the interval $[2, 10]$.
      Include appropriate units.

      \vspace{0.08in}
      $\text{ARC} = \dfrac{f(10) - f(2)}{10 - 2} = \dfrac{\blank{2.5cm}}{\blank{2cm}} = $ \blank{2.5cm}

      \vspace{0.04in}
      Units: \blank{4cm}

\item Evaluate $p(x) = -3x^3 + 2x^2 - x + 5$ at $x = -2$. Show each step.

      \vspace{0.08in}
      $p(-2) = -3(\blank{1cm})^3 + 2(\blank{1cm})^2 - (\blank{1cm}) + 5 = $ \blank{3cm}

\item A scientist says: ``As the distance between two magnets doubles, the force between them
      drops to one-fourth of its original value.'' Circle the model type that best fits
      this description and explain in one sentence.

      \medskip
      \textbf{A.}~Linear \qquad \textbf{B.}~Quadratic \qquad \textbf{C.}~Rational (inverse proportion) \qquad \textbf{D.}~Cubic

      \vspace{0.08in}
      Explanation: \writeline

\end{enumerate}

\end{document}
